\documentclass[11pt,a4paper]{article}

\usepackage[T1]{fontenc}
\usepackage[margin=2.2cm]{geometry}
\usepackage{listings}
\usepackage{xcolor}
\usepackage{booktabs}
\usepackage{enumitem}
\usepackage[colorlinks=true,linkcolor=blue!50!black,urlcolor=blue!50!black]{hyperref}

\lstset{
  basicstyle=\ttfamily\small,
  columns=fullflexible,
  keepspaces=true,
  frame=single,
  rulecolor=\color{gray!50},
  backgroundcolor=\color{gray!8},
  breaklines=true,
  showstringspaces=false,
  aboveskip=0.8em,
  belowskip=0.8em
}
\newcommand{\cmd}[1]{\texttt{#1}}
\setlist{itemsep=2pt,topsep=4pt}

\title{QEMU Debugging Reference for Kernel Development}
\author{}
\date{\today}

\begin{document}
\maketitle
\tableofcontents
\bigskip

\section{Choosing the Accelerator}
The most important decision is whether to run under KVM or under TCG
(QEMU's software emulator), because it determines which debugging
facilities exist.
\begin{itemize}
  \item The \cmd{-d} logging categories (\cmd{int}, \cmd{in\_asm},
        \cmd{exec}, \cmd{mmu}, \cmd{cpu\_reset}) work \textbf{only under TCG}.
        Under \cmd{-enable-kvm} the host kernel delivers exceptions and
        interrupts, so QEMU never observes them and \cmd{-d int} prints nothing.
  \item \cmd{-d guest\_errors} and \cmd{-d unimp} still work under KVM,
        because they log device-emulation events.
  \item TCG does not provide PCID/INVPCID or AVX-512, and its timing differs
        from real hardware. A feature check that passes under KVM can fail under TCG.
  \item Recommended workflow: reproduce the bug under KVM, then trace it under TCG.
\end{itemize}
\section{Command-Line Flags}
\begin{lstlisting}
qemu-system-x86_64 -accel tcg -cpu max -smp 1 \
  -m 512M -cdrom kernel.iso -display none \
  -serial mon:stdio \
  -no-reboot -no-shutdown \
  -d int,cpu_reset,guest_errors,unimp -D qemu.log \
  -debugcon file:debugcon.log \
  -device isa-debug-exit,iobase=0xf4,iosize=0x04
\end{lstlisting}

\begin{description}[style=nextline]
  \item[\cmd{-no-reboot -no-shutdown}]
    A triple fault freezes the VM instead of silently rebooting, so you can
    inspect it with the monitor.
  \item[\cmd{-d items -D file}]
    Enables logging and sends it to a file, so it does not interleave with
    serial output.
  \item[\cmd{-debugcon file:debugcon.log}]
    Captures bytes written with \cmd{outb(0xe9, c)}. It needs no driver, so it
    works before serial is initialised.
  \item[\cmd{-device isa-debug-exit}]
    \cmd{outb(0xf4, code)} makes QEMU exit with status \cmd{(code<<1)|1}, which
    is useful for automated tests.
  \item[\cmd{-smp 1}]
    Removes AP startup and SIPI traffic from the logs. Add CPUs back once the
    single-CPU path works.
  \item[\cmd{-cpu qemu64}]
    A strictness test. If your kernel accidentally depends on a host feature,
    it faults here but works under \cmd{-cpu host} or \cmd{max}.
  \item[\cmd{-s -S}]
    Starts a GDB server on TCP port 1234 and pauses at reset (see
    Section~\ref{sec:gdb}).
\end{description}

\subsection*{Additional \texttt{-d} items}

\begin{center}
\begin{tabular}{@{}ll@{}}
\toprule
Item & Meaning \\
\midrule
\cmd{int}         & Every exception and interrupt delivery, with register state \\
\cmd{cpu\_reset}  & CPU state dump on reset, including triple faults \\
\cmd{guest\_errors} & Invalid guest accesses to emulated devices \\
\cmd{unimp}       & Unimplemented device features the guest touched \\
\cmd{in\_asm}     & Guest code blocks as they are translated \\
\cmd{exec}        & Every executed translation block (very verbose) \\
\cmd{cpu}         & CPU state at each translation block \\
\cmd{mmu}         & MMU and page-table activity \\
\cmd{op}          & TCG operations (very verbose) \\
\bottomrule
\end{tabular}
\end{center}

Run \cmd{-d help} to list what your build supports. Items are combined with commas.
\section{Monitor Basics}

With \cmd{-serial mon:stdio}, the serial console and the QEMU monitor share the
terminal.

\begin{lstlisting}
Ctrl-A c       toggle between serial console and monitor
Ctrl-A x       quit QEMU

stop           pause the VM
cont           resume
system_reset   reset the guest
nmi            inject an NMI (needs a handler in your IDT)
help info      list `info` subcommands in your build
quit
\end{lstlisting}

Run \cmd{stop} before inspecting state, otherwise the registers change while you
read them. Command availability varies between QEMU versions, so check
\cmd{help info} on your build.

\section{CPU State and Disassembly}

\begin{lstlisting}
info registers          # current CPU
info registers -a       # all CPUs
info cpus               # per-CPU program counter and halted status

x/8i $rip               # disassemble 8 instructions at RIP
x/16gx $rsp             # 16 quadwords at RSP (virtual addresses)
p $rax                  # print a register
\end{lstlisting}

\subsection*{Reading an \texttt{info registers} dump}

\begin{itemize}
  \item \cmd{HLT=1} means the CPU executed \cmd{hlt}. With IF set, it is idle and
        waiting for an interrupt.
  \item QEMU prints RIP \emph{after} the \cmd{hlt}, so the instruction is
        usually at RIP$-1$.
  \item The flag string \cmd{[---Z-P-]} shows DF, OF, SF, ZF, AF, PF, CF only.
        IF is not shown, so read it from bit 9 of the \cmd{RFL} hex value.
  \item \cmd{CS64} in the CS line means the CPU is executing 64-bit code.
  \item \cmd{CPL} is the current privilege level.
\end{itemize}

\begin{center}
\begin{tabular}{@{}lll@{}}
\toprule
Register & Bit & Meaning \\
\midrule
CR0   & 0, 31, 16 & PE, PG, WP \\
CR4   & 4, 5, 9, 17, 18 & PSE, PAE, OSFXSR, PCIDE, OSXSAVE \\
EFER  & 0, 8, 10, 11 & SCE, LME, LMA, NXE \\
\bottomrule
\end{tabular}
\end{center}
\section{Memory and Paging}

\begin{lstlisting}
info mem                       # virtual mappings with permissions
info tlb                       # cached translations
gva2gpa 0xffffff8000101d46     # walk the page tables for one address
xp /8gx 0x109000               # read PHYSICAL memory (here: the PML4)
info mtree                     # physical address map
info mtree -f                  # flattened view

memsave  0xffffff8000100000 0x1000 virt.bin    # dump virtual memory
pmemsave 0x100000 0x1000 phys.bin              # dump physical memory
dump-guest-memory core.elf                     # full guest core file
\end{lstlisting}

To verify a mapping by hand, take CR3 from \cmd{info registers}, then read
PML4, PDPT, PD and PT entries with \cmd{xp} and compare them with what your
code wrote. \cmd{info mem} is the quickest check of a higher-half mapping.
\section{Interrupt Hardware}

\begin{lstlisting}
info pic        # 8259 IRR/IMR/ISR (IOAPIC pins in some builds)
info lapic      # LVT entries, timer, SPIV, ICR, IRR/ISR
info irq        # per-IRQ counts (not in every build)
info pci
info qtree      # device tree, e.g. confirm serial at 0x3f8
\end{lstlisting}

\begin{itemize}
  \item \cmd{imr=ff} on the master PIC: every legacy IRQ is masked.
  \item \cmd{irr} set and \cmd{isr} clear: the interrupt is pending but not delivered.
  \item \cmd{isr} stuck set: the handler never sent an EOI.
  \item LVT value \cmd{0x10000} has bit 16 set, meaning that entry is masked.
  \item A halted CPU with IF set and everything masked will never wake up.
\end{itemize}
\section{Reading the \texttt{-d int} Log}

\begin{lstlisting}
grep -nE 'check_exception|Triple fault|^ *[0-9]+: v=' qemu.log | tail -30
\end{lstlisting}

\begin{itemize}
  \item Each \cmd{v=XX e=XXXX ... CR2=...} line is one exception or interrupt
        delivery. \cmd{v} is the vector, \cmd{e} the error code.
  \item \cmd{check\_exception old: ... new: ...} shows the fault chain that ends
        in a double or triple fault.
  \item \cmd{cpu\_reset} dumps the full CPU state at the triple fault.
  \item Once timer interrupts are on, \cmd{int} logs every tick. Filter the log
        or drop \cmd{int} for long runs.
\end{itemize}
\section{GDB over QEMU}\label{sec:gdb}

\begin{lstlisting}
# terminal 1
qemu-system-x86_64 ... -s -S

# terminal 2
gdb -q boot.elf \
    -ex 'set architecture i386:x86-64' \
    -ex 'target remote :1234'
(gdb) hbreak kmain
(gdb) continue
(gdb) monitor info lapic       # any HMP command works
(gdb) info registers rip
\end{lstlisting}

\begin{itemize}
  \item Use \cmd{hbreak} instead of \cmd{break} in early boot. Software
        breakpoints patch memory through virtual addresses that may not be
        mapped yet.
  \item \cmd{watch} is more reliable under TCG than under KVM.
  \item GDB can get confused when the guest changes mode (real, protected,
        long). Re-run \cmd{set architecture i386:x86-64} if it does.
\end{itemize}
\section{Makefile Targets}
Appended to an existing Makefile, these add TCG debugging without changing
earlier rules. Recipe lines need tabs.
\begin{lstlisting}
DBG_ITEMS ?= int,cpu_reset,guest_errors,unimp
DBG_CPU   ?= max
DBG_SMP   ?= 1

QEMU_TCG_FLAGS := -m 512M -cdrom kernel.iso -serial mon:stdio \
                  -no-reboot -no-shutdown -display none \
                  -smp $(DBG_SMP) -accel tcg -cpu $(DBG_CPU) \
                  -d $(DBG_ITEMS) -D qemu.log \
                  -debugcon file:debugcon.log \
                  -device isa-debug-exit,iobase=0xf4,iosize=0x04

.PHONY: run-tcg run-tcg-gdb gdb
run-tcg: iso
	$(QEMU) $(QEMU_TCG_FLAGS)
run-tcg-gdb: iso
	$(QEMU) $(QEMU_TCG_FLAGS) -s -S
gdb: $(TARGET)
	gdb -q $(TARGET) -ex 'set architecture i386:x86-64' \
	    -ex 'target remote :1234'
\end{lstlisting}

\section{Troubleshooting Checklist}

\begin{center}
\small
\begin{tabular}{@{}p{0.34\linewidth}p{0.58\linewidth}@{}}
\toprule
Symptom & First things to try \\
\midrule
Reboot loop or silent hang
  & TCG with \cmd{-d int,cpu\_reset -no-reboot -no-shutdown}; read the tail of \cmd{qemu.log} \\
Silent after \cmd{sti; hlt}
  & \cmd{info pic} and \cmd{info lapic}: look for masks, pending bits, a missing IDT gate \\
Wrong or missing mapping
  & \cmd{info mem}, \cmd{gva2gpa}, then \cmd{xp} on the page tables \\
SMP oddities
  & Run with \cmd{-smp 1} first, then \cmd{info cpus} \\
CPU feature reported missing
  & Check whether you are running under TCG instead of KVM \\
Need source-level stepping
  & \cmd{-s -S} plus \cmd{hbreak} in GDB \\
\bottomrule
\end{tabular}
\end{center}

\end{document}
